\section{Conclusion}
\label{sec:conclusion}

We presented a controlled forward benchmark for neural emulation of solver-generated tsunami-like wavefields. The benchmark pairs procedurally generated bathymetry with diverse initial-displacement sources on shared scenarios, and labels each scenario with up to three numerical references -- a hydrostatic-reconstruction shallow-water solver as the primary reference, a second-order MUSCL-HR shallow-water solver, and an elevation-only weakly dispersive Boussinesq solver -- so that the same input scenario can be emulated and compared across solver fidelities without changing the underlying seabed or forcing. On this benchmark we studied FNO and FFNO surrogates against CNN, U-Net, U-FNO, WNO, and Fourier-mode ablations under a matched training and evaluation protocol, and we reported accuracy, runtime, distribution-shift behavior, cross-resolution transfer, synthetic-source transfer onto fully wet real-bathymetry morphologies, physics diagnostics, and predictive uncertainty.

The core findings are consistent across the diagnostics. On same-solver hydrostatic accuracy the FFNO is the strongest direct model (relative $L_2$ $0.168$), ahead of the base FNO and U-FNO near $0.235$, WNO at $0.367$, U-Net at $0.475$, and CNN at $0.607$. The base FNO emulates all three references to a comparable dimensionless fidelity ($0.21$--$0.26$), though we are careful to note that this ranking is unit-dependent and reverses in physical units. The single-pass models drift upward over the forecast horizon on the shallow-water references, and seeded windowed variants reduce that error substantially: Window-FNO-5 reaches $0.077$ and Window-FFNO-5 reaches $0.052$, at the cost of being seeded with the true first frame and running several sequential passes. Inference is several orders of magnitude faster than the CPU NumPy reference implementations used here, which is the regime change that makes large-ensemble and many-scenario studies practical. The strict held-out-family tests are mixed, ranging from mild trench degradation to a rough-source failure, and the 10-crop real-bathymetry suite supports only a modest fully wet morphology-transfer claim under synthetic forcing. The cross-solver diagnostic returns an honest negative result: the emulator-superiority ratio exceeds one in every pair, so a surrogate reproduces the solver it was trained on far better than it reproduces a different solver, and it does not close the physical gap between solvers. The deep ensemble gives an uncertainty signal that tracks error usefully (correlation $\approx 0.65$--$0.70$) but is systematically overconfident, with nominal intervals under-covering both in-distribution and under shift.

We emphasize that multi-reference evaluation is what makes these readings trustworthy rather than artifacts of a single label choice. Keeping the raw solver-vs-solver gap separate from the surrogate-vs-surrogate error prevents a surrogate's emulation skill from being confused with the physical disagreement between numerical models, and reporting every error against a stated solver and configuration keeps the claims honest: these are errors against numerical references, not against the real ocean.

Several directions follow naturally and are deliberately left to future work. Broader real-world transfer would require surveyed bathymetry, realistic rupture models, shoreline and wet-dry physics, and validation against observed records. The runtime comparison would be sharpened by benchmarking against optimized GPU or compiled finite-volume solvers rather than the reproducible CPU references used here. Stronger physical references, including a full nonlinear Boussinesq model, would let the cross-fidelity comparison speak to dispersive regimes more directly. The uncertainty estimates would benefit from recalibration and larger ensembles before any reliability claim. Finally, the inverse problem of source reconstruction from sparse observations -- for which the repository already contains a data scaffold -- needs different inputs, metrics and validation, and we leave it to a separate follow-up study rather than mixing it into the forward-surrogate benchmark.
